% SPDX-License-Identifier: CC-BY-4.0
% Copyright (C) 2026 Fabian Januszewski, Abhijit and the Langlandscape team
% Stage 03, WP S03-M1: weighted graphs, Phi^comb, the DR graph and the Mazur check.
% Labels are fixed by plan/staged/stage-03-ribet-ihara.tex. Owner of the interface if:s03-graph.
% Conventions (C1)--(C6) are in chapters/s03/4-ihara.tex (\S\ref{ss:s03-conv}); prop:s03-phieis lives in 2-eisenstein.tex (S03-M2).
\section{Weighted graphs and the Mazur check}\label{ss:s03-graphs}

\emph{Standing notation.} A \emph{graph} is a finite set $V$ of vertices, a finite set $E$ of edges and two maps
$o,t:E\to V$ (origin and target; loops $o(e)=t(e)$ and multiple edges are allowed). Its chain complex is
$\Z[E]\xrightarrow{\partial}\Z[V]$, $\partial e=t(e)-o(e)$, followed by $\deg:\Z[V]\to\Z$. The graph is \emph{connected} if $V\ne\emptyset$ and
the equivalence relation on $V$ generated by $o(e)\sim t(e)$ has one class (as in stage 02, \S\ref{ss:s02-quotient}). For a finite set $S$
we write $\Z[S]^0:=\ker(\deg:\Z[S]\to\Z)$, $\mathbf 1\in\Z^S$ for the constant function $1$, and $\mathbf 1_s$ for the indicator of $s\in S$.

\begin{definition}[weighted graph; seed \S2.7]\label{def:s03-graph}
A \emph{weighted graph} $\G$ is a finite graph (multiple edges and loops allowed) with vertex set $V(\G)$, edge set $E(\G)$, a
chosen orientation, and weights $w_e\in\Z_{\ge1}$. Its chain complex is $\Z[E]\xrightarrow\partial\Z[V]$. Put
$X(\G):=H_1(\G)=\ker\partial$, $\langle e,e'\rangle_w=w_e\delta_{ee'}$ and
\[\Phicomb(\G):=\coker\big(X(\G)\to\Hom(X(\G),\Z),\ x\mapsto\langle x,\cdot\rangle_w\big).\]
\status{definition} Lean (S03-L1): \lean{Langlandscape.Graph.WGraph}, \lean{Langlandscape.Graph.WGraph.boundary},
\lean{Langlandscape.Graph.WGraph.X} (alias \lean{Langlandscape.Graph.X}), \lean{Langlandscape.Graph.deg0},
\lean{Langlandscape.Graph.wPairing}, \lean{Langlandscape.Graph.wDualMap}, \lean{Langlandscape.Graph.phiComb},
\lean{Langlandscape.Graph.WGraph.PhiComb} \stformal{} (definitions; the Lean graph carries the positivity of the weights and has no finiteness
or connectedness field, finiteness is a hypothesis of the theorems).
\end{definition}

\noindent\emph{Basic facts} (\stproved). (1) $X(\G)$ is a direct summand of $\Z[E]$ (the quotient $\Z[E]/X(\G)\cong\partial\Z[E]\subseteq\Z[V]$
is free), so it is free of finite rank, and restriction $\Hom(\Z[E],\Z)=\Z^E\to\Hom(X(\G),\Z)$ is surjective. (2) The form
$\langle,\rangle_w$ is positive definite on $\Z[E]\otimes\R$, hence on $X(\G)\otimes\R$, so $x\mapsto\langle x,\cdot\rangle_w$ is injective. If
$b_1,\dots,b_g$ is a basis of $X(\G)$ with Gram matrix $G=(\langle b_i,b_j\rangle_w)$, then in the bases $(b_i)$ and the dual basis the map is
given by $G$, and
\[|\Phicomb(\G)|=\det G>0 .\]
Indeed the cokernel of an injective endomorphism of $\Z^g$ with matrix $G$ has order $|\det G|$ (Smith normal form), and $\det G>0$ for a
positive definite form. (3) Reversing the orientation of an edge $e$ replaces the basis vector $e$ by $-e$. This is an isometry of
$(\Z[E],\langle,\rangle_w)$ which carries the old $X(\G)$ onto the new one, so $\Phicomb(\G)$ does not depend on the orientation up to
isomorphism.

\begin{remark}\label{rem:s03-grothendieck}
For a semistable curve with dual graph $\G$ and thicknesses $w_e$, $X(\G)$ is the character group of the toric part. Then
$\langle,\rangle_w$ is the monodromy pairing and $\Phicomb$ is the component group: [SGA7, IX] \stcited{} (reference not in
\pathtt{literature/}), and [Ri90, Prop.~2.2, Thm~2.3, p.~436; weighted form p.~437--438]
\lit{ribet\_1990\_modular-representations-of-gal-qbar-q-arising-from-modular-forms}{Prop.~2.2, Thm.~2.3, p.~436; weighted form
pp.~437--438}. The tree route does not use this dictionary. In this plan it is a remark only, and its use belongs to stage 05.
\end{remark}

\begin{proposition}[two-vertex graphs; seed Prop.~phi]\label{prop:s03-phi}
Let $\G$ have two vertices $v_0\ne v_1$ and $E\ne\emptyset$, and let every edge be oriented from $v_0$ to $v_1$. Then $X(\G)=\Z[E]^0$ and $\Phicomb(\G)=\Z^E/(W\Z[E]^0+\Z\mathbf1)$, where
$W=\mathrm{diag}(w_e)$. Moreover
\[|\Phicomb(\G)|=\Big(\prod_ew_e\Big)\sum_e\frac1{w_e}=\sum_e\prod_{e'\ne e}w_{e'}.\]
If all $w_e=1$, then $\Phicomb\cong\Z/|E|$. In general $\Phicomb$ is not cyclic: $w=(2,2,2)$ gives $\Z/2\times\Z/6$.
For an arbitrary orientation of the edges the order formula and the isomorphism class of $\Phicomb$ are unchanged (fact~(3)), but $X(\G)$ is
$\Z[E]^0$ only after the sign change of the reversed edges. For $E=\emptyset$ one has $\Phicomb=0$ of order $1$ while the sum is $0$, so the
formula fails; this case is excluded.
\stproved{} (proof below; finite linear algebra). Lean (S03-L3): \lean{Langlandscape.Graph.twoVertex\_X\_eq\_deg0},
\lean{Langlandscape.Graph.phiComb\_twoVertex\_equiv}, \lean{Langlandscape.Graph.card\_phiComb\_twoVertex},
\lean{Langlandscape.Graph.phiComb\_twoVertex\_unit} \stformal{}, for $E\ne\emptyset$ and all edges oriented $v_0\to v_1$ (the Lean hypotheses
\texttt{Nonempty G.E}, \texttt{hsrc}, \texttt{htgt}). The orientation-independence step (fact~(3)) and the example $w=(2,2,2)$ are not formalized.
\end{proposition}
\begin{proof}
By hypothesis $E\ne\emptyset$ (for $E=\emptyset$, $\Phicomb=0$ has order $1$ and the sum is $0$, so the formula fails; this case is excluded and does not
occur below) and every edge is oriented from $v_0$ to $v_1$ (for another orientation one reduces to this case by fact~(3), for the order and the
isomorphism class). Then $\partial\big(\sum_ex_ee\big)=\big(\sum_ex_e\big)(v_1-v_0)$,
so $X(\G)=\Z[E]^0$.

\emph{The presentation.} $\Z[E]=\Z[E]^0\oplus\Z e_*$ for any fixed edge $e_*$, so restriction $\Z^E\to\Hom(\Z[E]^0,\Z)$ is onto, with kernel
the functionals vanishing on $\Z[E]^0$, i.e.\ $\Z\mathbf 1$ (a functional $\varphi$ kills all $e-e_*$ iff it is constant). Hence
$\Hom(X(\G),\Z)=\Z^E/\Z\mathbf 1$. Under this identification $x\in\Z[E]^0$ goes to the class of $e\mapsto\langle x,e\rangle_w=w_ex_e$, i.e.\ of
$Wx$. So $\Phicomb(\G)=\Z^E/(W\Z[E]^0+\Z\mathbf 1)$.

\emph{The order.} Number the edges $e_1,\dots,e_n$ and use the basis $b_i=e_i-e_n$ ($1\le i\le n-1$) of $\Z[E]^0$. Its Gram matrix is
$G=\langle b_i,b_j\rangle_w=w_i\delta_{ij}+w_n$, i.e.\ $G=D+w_n\mathbf u\mathbf u^{\mathsf T}$ with $D=\mathrm{diag}(w_1,\dots,w_{n-1})$ and
$\mathbf u=(1,\dots,1)^{\mathsf T}$. By the matrix determinant lemma $\det(D+c\,\mathbf u\mathbf u^{\mathsf T})=\det D\,(1+c\,\mathbf u^{\mathsf T}D^{-1}\mathbf u)$,
\[\det G=\prod_{i<n}w_i\Big(1+w_n\sum_{i<n}\frac1{w_i}\Big)=\prod_{i\le n}w_i\sum_{i\le n}\frac1{w_i}=\sum_e\prod_{e'\ne e}w_{e'},\]
and $|\Phicomb(\G)|=\det G$ by fact (2). (For $n=1$, $G$ is the empty matrix, $\det G=1$, and indeed $\Phicomb=0$ and $\sum_e\prod_{e'\ne e}w_{e'}=1$.)

\emph{All $w_e=1$.} Then $\gcd_e(W/w_e)=1$ in Lemma~\ref{lem:s03-cyclic} (each $W/w_e=1$), so $\Phicomb\cong\Z/n$ with $n=|E|$.

\emph{$w=(2,2,2)$.} With $b_1,b_2$ as above, $G=\left(\begin{smallmatrix}4&2\\2&4\end{smallmatrix}\right)$. The gcd of its entries is $2$ and
$\det G=12$, so its Smith normal form is $\mathrm{diag}(2,6)$ and $\Phicomb\cong\Z/2\times\Z/6$, which is not cyclic. (Lemma~\ref{lem:s03-cyclic}
agrees: $\gcd_e(W/w_e)=\gcd(4,4,4)=4\ne1$.)
\end{proof}

\begin{lemma}[cyclicity]\label{lem:s03-cyclic}
In the situation of Prop.~\ref{prop:s03-phi}, put $W:=\prod_ew_e$, $n:=|\Phicomb(\G)|$ and $\lambda(x):=\sum_e(W/w_e)x_e$. Then
$\lambda$ induces a homomorphism $\bar\lambda:\Phicomb(\G)\to\Z/n$. It is surjective, hence an isomorphism, if and only if
$\gcd_e(W/w_e)=1$. This holds in particular when the $w_e$ are pairwise coprime (any number of $1$'s allowed). \stproved{} (proof below;
the stage tex has \sttoprove). Lean (S03-L4): \lean{Langlandscape.Graph.lambdaBar},
\lean{Langlandscape.Graph.lambdaBar\_surjective\_iff}, \lean{Langlandscape.Graph.lambdaBar\_bijective\_of\_pairwise\_coprime} \stformal{}
(for two-vertex graphs with $E\ne\emptyset$, edges $v_0\to v_1$), with $\lambda$ itself as \lean{Langlandscape.Graph.lambdaMap} \stformal{} (a Lean definition on $\Z[E]$ with the
same formula).
\end{lemma}
\begin{proof}
In this lemma $W$ is the integer $\prod_ew_e$; the diagonal matrix of Prop.~\ref{prop:s03-phi} is written $\mathrm{diag}(w_e)$. Use the presentation
$\Phicomb(\G)=\Z^E/(\mathrm{diag}(w_e)\Z[E]^0+\Z\mathbf 1)$ of Prop.~\ref{prop:s03-phi}, and regard $\lambda$ as a homomorphism $\Z^E\to\Z$.

\emph{Well defined.} $\lambda(\mathbf 1)=\sum_eW/w_e=W\sum_ew_e^{-1}=n$ by Prop.~\ref{prop:s03-phi}, and for $x\in\Z[E]^0$,
$\lambda(\mathrm{diag}(w_e)x)=\sum_e(W/w_e)w_ex_e=W\sum_ex_e=0$. So $\lambda$ maps the denominator into $n\Z$ and induces $\bar\lambda$.

\emph{Image.} $\lambda(\Z^E)=d\Z$ with $d:=\gcd_e(W/w_e)$ (the image of $\mathbf 1_e$ is $W/w_e$). Since $n=\sum_eW/w_e$, $d\mid n$, so the image of
$\bar\lambda$ is $d\Z/n\Z$, a group of order $n/d$. Hence $\bar\lambda$ is onto iff $d=1$. If it is onto, it is a surjection between finite groups of the same
order $n$, hence bijective.

\emph{Pairwise coprime weights.} Suppose a prime $r$ divides $d$. Then $r\mid W/w_e$ for every $e$; in particular $r\mid W$, so $r\mid w_{e_0}$ for some
$e_0$. As the weights are pairwise coprime, $r\nmid w_e$ for $e\ne e_0$, so $r\nmid\prod_{e\ne e_0}w_e=W/w_{e_0}$, a contradiction. Hence $d=1$.
\end{proof}

\begin{remark}\label{rem:s03-kirchhoff}
For a connected weighted graph, $|\Phicomb(\G)|=\sum_{T}\prod_{e\notin T}w_e$, where $T$ runs over the spanning trees. This is the
weighted matrix-tree theorem. \stproved{} (proof below; the stage tex has \sttoprove{} (standard)). Prop.~\ref{prop:s03-phi} is the two-vertex case:
there the spanning trees are the single edges.
\end{remark}
\begin{proof}
Let $\G$ be connected, $|V|=m\ge1$, $g:=|E|-m+1=\rk X(\G)$ (Euler characteristic of the exact sequence
$0\to X(\G)\to\Z[E]\to\Z[V]\xrightarrow{\deg}\Z\to0$; exactness at $\Z[V]$ is connectedness). Fix a spanning tree $T_0$. For an edge
$f\notin T_0$ let $c_f\in X(\G)$ be its fundamental cycle ($f$ plus the unique path in $T_0$ from $t(f)$ to $o(f)$). The $c_f$ form a basis of
$X(\G)$: for $x\in X(\G)$, $x-\sum_{f\notin T_0}x_fc_f$ is a cycle supported on $T_0$, hence $0$ because $H_1(T_0)=0$. Let $C$ be the $g\times|E|$
matrix of the $c_f$ in the basis $E$. The Gram matrix is $G=C\,\mathrm{diag}(w_e)\,C^{\mathsf T}$, so by the Cauchy--Binet formula
\[\det G=\sum_{S\subseteq E,\ |S|=g}\det(C_S)^2\prod_{e\in S}w_e,\]
where $C_S$ is the $g\times g$ minor on the columns $S$. Now $C_S$ is the matrix of the restriction $\rho_S:X(\G)\to\Z^S$,
$x\mapsto(x_e)_{e\in S}$. Its kernel is $X(\G)\cap\Z[E\setminus S]=H_1$ of the subgraph with edge set $T:=E\setminus S$ (and all vertices). This subgraph has
$m-1$ edges. If $T$ is a spanning tree, then $\rho_S$ is an isomorphism (surjective because the fundamental cycles $c^T_f$, $f\in S$, of $T$ map to the
standard basis of $\Z^S$; injective as $H_1(T)=0$), so $\det C_S=\pm1$. If $T$ is not a spanning tree, then the graph $(V,T)$, which has $m$ vertices and
$m-1$ edges, is not a tree, hence not acyclic, so $H_1(V,T)\ne0$, $\rho_S$ is not injective and $\det C_S=0$. Hence
$\det G=\sum_T\prod_{e\notin T}w_e$, and $|\Phicomb(\G)|=\det G$ by fact (2) after Def.~\ref{def:s03-graph}.
\end{proof}

\begin{definition}[DR graph; seed Def.~DR]\label{def:s03-DR}
Let $B=B_q$ and $U=\Uone{Mp^s}$ or $\Uzero M$. Then $\GDR_q(U)$ is the two-vertex weighted graph with edge set $\XBq U$, every edge
oriented $0\to1$, and $w_e=|\Gamma_e/(\Gamma_e\cap\{\pm1\})|$. Put $\Xcomb_q(U):=\Z[\XBq U]^0$ and, for $k>2$,
$\Xcomb_q(U;\Vk):=\Brandt U{\Vk}$. \status{definition} Lean (S03-L6): \lean{Langlandscape.Ribet.DRgraph} at
$U=\hat\cO^\times$ \stformal{}
\end{definition}

For $s\ge1$ every $w_e=1$ by Lemma~\ref{lem:s01-stab}(b) (\pathtt{if:s01-classset}; here $p\ge5$, $p\ne q$), so
$\Phicomb(\GDR_q(U))\cong\Z/|S_s|$ by Prop.~\ref{prop:s03-phi}, with $|S_s|=|\XBq{\Uone{Mp^s}}|$ given by Thm.~\ref{thm:s01-mass}(d). \stproved
By Prop.~\ref{prop:s03-phi}, $X(\GDR_q(U))=\Xcomb_q(U)$.

\subsection{The Mazur check}

\begin{lemma}[elliptic classes]\label{lem:s03-elliptic}
Let $q\ge5$, $B=B_q$ and $U=\hat\cO^\times$ maximal. Then $w_x\in\{1,2,3\}$ for every $[x]\in\XBq U$. Moreover
\[n_2:=\#\{x:w_x=2\}=\tfrac12\big(1-(\tfrac{-4}q)\big),\qquad n_3:=\#\{x:w_x=3\}=\tfrac12\big(1-(\tfrac{-3}q)\big).\]
That is, $n_2=1$ if and only if $q\equiv3\bmod4$, and $n_3=1$ if and only if $q\equiv2\bmod3$; otherwise each is $0$.
\stpmc: the mass formula \texttt{if:s01-mass} (Thm.~\ref{thm:s01-mass}(a)) [Vo, Thm~25.1.1, p.~415]
\lit{voight\_2021\_quaternion-algebras}{Thm 25.1.1, p.~415} and the Eichler class number formula [Vo, Thm~30.1.5, p.~543]
\lit{voight\_2021\_quaternion-algebras}{Thm 30.1.5, p.~543}, both \stcited; the list of unit groups [Vo, Thm~11.5.14, p.~176]
\lit{voight\_2021\_quaternion-algebras}{Thm 11.5.14, p.~176} \stcited.
Geometric form: [Mz, App.\ Thm~1.1(d), p.~176] \lit{mazur\_1977\_modular-curves-and-the-eisenstein-ideal}{App.\ Thm 1.1(d), p.~176}.
Lean: the statement is the registered hypothesis \pathtt{Langlandscape.Hyp.EllipticCount} (S03-L6), which also contains $|X_B(\hat\cO^\times)|=1$ for $q\in\{2,3\}$; the chapter proves that case from the mass formula and $|\cO_x^\times|\le24$ (Prop.~\ref{prop:s03-mazur}), so the hypothesis could be weakened to $q\ge5$; in Lean this is the registered hypothesis, not a theorem; \lean{Langlandscape.Ribet.mazur\_check} uses its $q\le3$ clause
\end{lemma}
\begin{proof}
\emph{Step 1: the possible weights.} For $U=\hat\cO^\times$, $\Gamma_x=\cO_x^\times$ with $\cO_x=B\cap x\hat\cO x^{-1}$ a maximal order, and
$|\cO_x^\times|\in\{2,4,6,12,24\}$ (Prop.~\ref{prop:s01-finite}). By [Vo, Thm~11.5.14] $\cO_x^\times$ is cyclic of order $2,4,6$, or $Q_8$, or binary
dihedral of order $12$, or binary tetrahedral of order $24$; the case $Q_8$ (the Lipschitz order) is excluded since $\cO_x$ is maximal
(Prop.~\ref{prop:s01-finite}); independently, $Q_8$ lies in the tetrahedral group, whose case below gives $B\cong(-1,-1)_\Q$. We show that the two non-cyclic cases force $B\not\cong B_q$ for $q\ge5$.
\begin{itemize}
\item \emph{Binary tetrahedral.} It contains a copy of $Q_8=\{\pm1,\pm i,\pm j,\pm ij\}$ with $i^2=j^2=-1$, $ij=-ji$. Then $1,i,j,ij$ span a
  $4$-dimensional subalgebra (the quaternion relations make them linearly independent over $\Q$), so $B\cong(-1,-1)_\Q$.
\item \emph{Binary dihedral of order $12$.} It contains an element $\omega$ of order $6$ and an element $j$ with $j^2=-1$ and
  $j\omega j^{-1}=\omega^{-1}$. Then $\omega^2-\omega+1=0$ (as $\omega$ has order $6$ in a division algebra, its minimal polynomial is the
  $6$th cyclotomic polynomial), so $s:=2\omega-1$ satisfies $s^2=4\omega^2-4\omega+1=-3$, and $jsj^{-1}=2\omega^{-1}-1=2(1-\omega)-1=-s$. Hence
  $B\cong(-3,-1)_\Q$.
\end{itemize}
Ramification via Hilbert symbols (for $\ell$ odd, a unit $u$ and $v(a)=1$: $(a,u)_\ell=(\tfrac u\ell)$; units pair trivially at odd $\ell$).
$(-1,-1)$: trivial at every odd $\ell$, non-trivial at $\infty$, hence (product formula) non-trivial at $2$: ramified exactly at
$\{2,\infty\}$. $(-3,-1)$: trivial at $\ell\nmid6$, non-trivial at $\infty$, $(-3,-1)_3=(\tfrac{-1}3)=-1$, hence trivial at $2$:
ramified exactly at $\{3,\infty\}$. Since $B_q$ is ramified exactly at $\{q,\infty\}$ and $q\ge5$, neither case occurs. So $\cO_x^\times$ is cyclic of
order $2,4,6$, and $w_x=|\cO_x^\times/\{\pm1\}|\in\{1,2,3\}$.

\emph{Step 2: the counts.} Put $h:=|\XBq U|$, $\epsilon_2:=1-(\tfrac{-4}q)$, $\epsilon_3:=1-(\tfrac{-3}q)\in\{0,2\}$. The mass formula gives
$\sum_xw_x^{-1}=\frac{q-1}{12}$ (Thm.~\ref{thm:s01-mass}(a) with $\varphi(q)=q-1$), and the Eichler class number formula for $D=q$ and maximal
level gives $h=\frac{q-1}{12}+\frac{\epsilon_2}4+\frac{\epsilon_3}3$ (as in Example~\ref{ex:s01-masschecks}). Subtracting, and using Step~1,
\[\sum_x\big(1-w_x^{-1}\big)=\frac{n_2}2+\frac{2n_3}3=\frac{\epsilon_2}4+\frac{\epsilon_3}3 .\]
Multiplying by $6$: $3n_2+4n_3=3\frac{\epsilon_2}2+4\frac{\epsilon_3}2\in\{0,3,4,7\}$. In non-negative integers, $3n_2+4n_3=0,3,4,7$ has the unique
solutions $(0,0),(1,0),(0,1),(1,1)$ respectively (if $n_3\ge2$ the left side is $\ge8$; for $n_3\in\{0,1\}$, $n_2$ is determined). Hence
$n_j=\epsilon_j/2$. Finally $(\tfrac{-4}q)=(\tfrac{-1}q)=-1$ iff $q\equiv3\bmod4$, and $(\tfrac{-3}q)=(\tfrac q3)=-1$ iff $q\equiv2\bmod3$ (quadratic
reciprocity).
\end{proof}

\begin{proposition}[Mazur check; seed Example mazur]\label{prop:s03-mazur}
For every prime $q$,
\[\Phicomb\big(\GDR_q(\hat\cO^\times)\big)\ \cong\ \Z\big/\mathrm{num}\tfrac{q-1}{12},\]
and this group is cyclic, generated via $\bar\lambda$ of Lemma~\ref{lem:s03-cyclic}. For $q\in\{2,3\}$ both sides are trivial.
\stpmc: Lemma~\ref{lem:s03-elliptic}, Prop.~\ref{prop:s03-phi}, Lemma~\ref{lem:s03-cyclic}. Lean (S03-L5, L6):
\lean{Langlandscape.Graph.mazurNum}, \lean{Langlandscape.Graph.mazur\_prod\_mul\_sum}, \lean{Langlandscape.Graph.mazur\_weights\_pairwise\_coprime}
\stformal{} (abstract arithmetic core: for weights in $\{1,2,3\}$ with the counts $n_2,n_3$ and $\sum w^{-1}=\frac{q-1}{12}$ one has
$\prod w\cdot\sum w^{-1}=\mathrm{num}\frac{q-1}{12}$, and the weights are pairwise coprime; $q\ge5$ only); the identification of the weights with the DR graph
(\lean{Langlandscape.Ribet.sum\_inv\_wx\_Ohat}, \lean{Langlandscape.Ribet.mazur\_check}, \stformal{} modulo the registered hypothesis \lean{Langlandscape.Hyp.EllipticCount} (and \pathtt{Hyp.EichlerMass})); the cases $q=2,3$ are \stformal{} as well, through the $q\le3$ clause of the same hypothesis \lean{Langlandscape.Hyp.EllipticCount} (which asserts $|\XB{\hat\cO^\times}|=1$ there; \lean{Langlandscape.Ribet.mazur\_check} covers every prime $q$), so the chapter's elementary derivation of that case is not formalized
\end{proposition}
\begin{proof}
Let $q\ge5$ and write $E=\XBq{\hat\cO^\times}$. By the mass formula, $\sum_e w_e^{-1}=\frac{q-1}{12}$. By Lemma~\ref{lem:s03-elliptic},
$\prod_ew_e=2^{n_2}3^{n_3}$, and $2^{n_2}3^{n_3}=12/\gcd(q-1,12)$. Indeed, $q$ is odd, so $\gcd(q-1,4)$ is $4$ if $q\equiv1\bmod4$ ($n_2=0$) and
$2$ if $q\equiv3\bmod4$ ($n_2=1$); and $\gcd(q-1,3)$ is $3$ if $q\equiv1\bmod3$ ($n_3=0$) and $1$ if $q\equiv2\bmod3$ ($n_3=1$). So with
$g:=\gcd(q-1,12)$, Prop.~\ref{prop:s03-phi} gives
\[|\Phicomb|=\prod_ew_e\sum_ew_e^{-1}=\frac{12}g\cdot\frac{q-1}{12}=\frac{q-1}g=\mathrm{num}\frac{q-1}{12}.\]
The weights lie in $\{1,2,3\}$, and $2$ and $3$ each occur at most once ($n_2,n_3\le1$), so the weights are pairwise coprime and
Lemma~\ref{lem:s03-cyclic} gives $\bar\lambda:\Phicomb\xrightarrow\sim\Z/\mathrm{num}\frac{q-1}{12}$; a generator is $\bar\lambda^{-1}(1)$.

Let $q\in\{2,3\}$. The mass is $\sum_xw_x^{-1}=\frac{q-1}{12}\in\{\frac1{12},\frac16\}$. Every $w_x\le12$ (as $|\cO_x^\times|\le24$), and for $q=3$
even $w_x\le6$: the order $24$ forces $B\cong(-1,-1)_\Q$ (Step~1 of the proof of Lemma~\ref{lem:s03-elliptic}), which is ramified at $2$, not at $3$.
So each term is $\ge\frac1{12}$ resp.\ $\ge\frac16$, and there is exactly one class, with $w=12$ resp.\ $6$. A two-vertex graph with one edge has
$X(\G)=\Z[E]^0=0$ and $\Phicomb=0$, and $\mathrm{num}\frac1{12}=\mathrm{num}\frac16=1$.
\end{proof}

\begin{remark}[comparison with Mazur; seed doubts]\label{rem:s03-mazur}
(a) Mazur--Rapoport--Deligne show that the component group of $J_0(q)$ at $q$ is cyclic of order $\mathrm{num}\frac{q-1}{12}$,
generated by $(0)-(\infty)$: [Mz, App.\ Thm~A.1(b), p.~173]
\lit{mazur\_1977\_modular-curves-and-the-eisenstein-ideal}{App.\ Thm A.1(b), p.~173, PDF p.~142} \stcited. The definite
statement above agrees in order \emph{and} in group structure. The identification of the two groups uses the dictionary of
Remark~\ref{rem:s03-grothendieck}: edges are supersingular points [Ri90, Prop.~3.1, Prop.~3.3, p.~439--440]
\lit{ribet\_1990\_modular-representations-of-gal-qbar-q-arising-from-modular-forms}{Prop.~3.1, Prop.~3.3, pp.~439--440}, and the
thickness is $e(i)=\#\mathrm{Aut}(E)/2=w_e$ [Ri90, p.~439]
\lit{ribet\_1990\_modular-representations-of-gal-qbar-q-arising-from-modular-forms}{p.~439},
[DR, VI Thm~6.9(iii), p.~286--287]
\lit{deligne-rapoport\_1973\_les-schemas-de-modules-de-courbes-elliptiques}{VI Th.~6.9(iii), pp.~286--287} \stcited. This identification is
\emph{not} part of the tree route.
(b) The seed restricts to $q\ge5$, but the statement holds for all $q$ (Prop.~\ref{prop:s03-mazur}).
(c) The seed's mass formula ``$\sum_e|\cO_e^\times|^{-1}=\frac{q-1}{24}$'' is correct with full unit groups (Thm.~\ref{thm:s01-mass}(a)). With
$w_e$ it reads $\frac{q-1}{12}$.
\end{remark}

\begin{example}[T1 data]\label{ex:s03-mazur}
Computed by HACS/PARI (WP S03-C2): the Brandt data of $B_q$ at level $\hat\cO^\times$ with $w_e=|\Gamma_e|/2$, and the Smith form of the Gram matrix of $X(\GDR_q)$.
\stcomp{}, \pathtt{code/results/s03/T1.json} (provenance: commit \texttt{518e93aa}, script \pathtt{code/python/langlandscape/ribet/mazur.py}, PARI 2.15.4).
\begin{center}\small
\begin{tabular}{rrlrrll}\toprule
$q$ & $h=|\XBq1|$ & weights $w_e$ & $\sum_ew_e^{-1}$ & $|\Phicomb|$ & $\Phicomb$ & Eisenstein\\\midrule
$11$ & $2$ & $(2,3)$ & $5/6$ & $5$ & $\Z/5$ & yes\\
$23$ & $3$ & $(1,2,3)$ & $11/6$ & $11$ & $\Z/11$ & yes\\
$37$ & $3$ & $(1,1,1)$ & $3$ & $3$ & $\Z/3$ & yes\\
$101$ & $9$ & $(1^8,3)$ & $25/3$ & $25$ & $\Z/25$ & yes\\
$131$ & $12$ & $(1^{10},2,3)$ & $65/6$ & $65$ & $\Z/65$ & yes\\\bottomrule
\end{tabular}
\end{center}
In every row $\Phicomb$ is cyclic (the Smith form has a single invariant factor $>1$), $|\Phicomb|$ equals $\mathrm{num}\frac{q-1}{12}$, and $\sum_ew_e^{-1}=\frac{q-1}{12}$
(mass formula); the numbers of the file agree with the hand values below, no discrepancy. $|\Phicomb|=\mathrm{num}\frac{q-1}{12}$ in each row. The last column records that $T_\ell-(1+\ell)$ kills $\Phicomb$, the check of
Prop.~\ref{prop:s03-phieis} for $\ell\in\{2,3,5,7,11,13\}\setminus\{q\}$ (\pathtt{eisenstein\_all} is true in every row). The class numbers are $(q-1)/12$, $(q+7)/12$, $(q+5)/12$, $(q+13)/12$ for
$q\equiv1,5,7,11\bmod12$; this is the formula of [Vo, Thm~30.1.5] \lit{voight\_2021\_quaternion-algebras}{Thm 30.1.5, p.~543} evaluated in the four residue
classes; they are listed as \lit{voight\_2021\_quaternion-algebras}{Exercise 30.6, p.~566}.
\end{example}
\noindent\emph{Hand check of the table.} For each $q$: $n_2=1$ iff $q\equiv3\bmod4$ and $n_3=1$ iff $q\equiv2\bmod3$ (Lemma~\ref{lem:s03-elliptic}), and the
remaining $h-n_2-n_3$ weights are $1$. E.g.\ $q=101\equiv5\bmod12$: $h=108/12=9$, $n_2=0$, $n_3=1$, $\sum w^{-1}=8+\frac13=\frac{100}{12}$, and
$\mathrm{num}\frac{100}{12}=25$; $q=131\equiv11\bmod12$: $h=144/12=12$, $n_2=n_3=1$, $\sum w^{-1}=10+\frac56=\frac{130}{12}$, $\mathrm{num}=65$. In each row
$\sum w^{-1}=\frac{q-1}{12}$, as the mass formula requires.

\begin{interface}[\pathtt{if:s03-graph}]\label{if:s03-graph}
Weighted graphs, $X(\G)$, $\langle,\rangle_w$, $\Phicomb$ (Def.~\ref{def:s03-graph}); Prop.~\ref{prop:s03-phi} and
Lemma~\ref{lem:s03-cyclic}; the DR graph $\GDR_q(U)$ and $\Xcomb_q(U)$ (Def.~\ref{def:s03-DR}); the Mazur check
Prop.~\ref{prop:s03-mazur} ($\Phicomb(\GDR_q(1))\cong\Z/\mathrm{num}\frac{q-1}{12}$, cyclic, all $q$); $\Phicomb$ is Eisenstein
(Prop.~\ref{prop:s03-phieis}, \S2 of this chapter, WP S03-M2). Status: \stpmc{} (Prop.~\ref{prop:s03-phi}, Lemma~\ref{lem:s03-cyclic},
Remark~\ref{rem:s03-kirchhoff} \stproved; Lemma~\ref{lem:s03-elliptic} and Prop.~\ref{prop:s03-mazur} \stpmc{} on the mass and class number
formulas). Lean names: S03-L1, L3--L6 rows of \pathtt{plan/implementation/stage-03.md}.
\end{interface}
